\section{Basic linear algebra}

Let $V$ be a vector space over $\mathbb{R}$ or $\mathbb{C}$. Vectors $\mathbf{v}_1,\ldots,\mathbf{v}_n$ are linearly independent when
\begin{equation}
c_1\mathbf{v}_1+\cdots+c_n\mathbf{v}_n=\mathbf{0}
\quad\Rightarrow\quad c_1=\cdots=c_n=0.
\end{equation}
If every vector in $V$ has a unique representation as a linear combination of these vectors, they form a basis of $V$ and $n$ is its dimension.

An $m\times n$ matrix $A=(a_{ij})$ maps vectors in $\mathbb{R}^n$ to vectors in $\mathbb{R}^m$. For square matrices,
\begin{equation}
\det(AB)=\det A\det B,
\end{equation}
and $A$ is invertible exactly when $\det A\ne0$. In low dimensions,
\begin{equation}
\det\begin{pmatrix}a&b\\c&d\end{pmatrix}=ad-bc,
\end{equation}
and
\begin{equation}
\det\begin{pmatrix}a_1&a_2&a_3\\b_1&b_2&b_3\\c_1&c_2&c_3\end{pmatrix}
=a_1(b_2c_3-b_3c_2)-a_2(b_1c_3-b_3c_1)+a_3(b_1c_2-b_2c_1).
\end{equation}

The standard inner product on $\mathbb{R}^n$ is
\begin{equation}
\langle\mathbf{u},\mathbf{v}\rangle=\mathbf{u}\cdot\mathbf{v}=\sum_i u_iv_i.
\end{equation}
Vectors are orthogonal when their inner product vanishes. More generally, a function-space inner product may be defined by
\begin{equation}
\langle f,g\rangle=\int_a^b f(x)g(x)\,\mathrm{d}x,
\end{equation}
which is the foundation of Fourier orthogonality.

An eigenvalue $\lambda$ and eigenvector $\mathbf{v}\ne\mathbf{0}$ of a square matrix $A$ obey
\begin{equation}
A\mathbf{v}=\lambda\mathbf{v}.
\end{equation}
They are found from the characteristic equation $\det(A-\lambda I)=0$. A real symmetric matrix has real eigenvalues, and eigenvectors belonging to distinct eigenvalues are orthogonal. The linear system $A\mathbf{x}=\mathbf{b}$ has the unique solution $\mathbf{x}=A^{-1}\mathbf{b}$ when $\det A\ne0$.

Given linearly independent vectors $f_1,\ldots,f_n$, Gram--Schmidt orthogonalization constructs
\begin{equation}
g_1=f_1,
\qquad
g_k=f_k-\sum_{j=1}^{k-1}\frac{\langle f_k,g_j\rangle}{\langle g_j,g_j\rangle}g_j.
\end{equation}
Normalizing the $g_k$ produces an orthonormal basis, a construction that reappears in expansions by orthogonal functions.
